\documentclass{article}%
\usepackage[T1]{fontenc}%
\usepackage[utf8]{inputenc}%
\usepackage{lmodern}%
\usepackage{textcomp}%
\usepackage{lastpage}%
\usepackage{geometry}%
\geometry{margin=2.2cm}%
%
\usepackage{amsmath}%
\usepackage{amssymb}%
\usepackage{booktabs}%
\usepackage{longtable}%
\usepackage{array}%
\usepackage{hyperref}%
\hypersetup{hidelinks}%
\title{ITA 2027 {-} Matemática (claude\_haiku\_5\_5) x Poliedro\\[4pt]\large Comparação com o gabarito Poliedro Resolve (ITA 2027, 2ª fase, dia 1 {-} Matemática)\thanks{Gerado com gpt-resolve: \url{https://github.com/lgabs/gpt-resolve}}}%
\date{09/11/2026}%
%
\begin{document}%
\normalsize%
\maketitle%
\section*{Resumo}\textbf{Nota média: 10,0/10} (12 de 12 itens com a resolução correta; cada item vale o mesmo dentro da questão).\par\medskip%
\begin{longtable}{@{}c c p{0.78\textwidth}@{}}\toprule Questão & Nota & Resumo \\ \midrule
1 & 10,0 & A resolução da Q1 está correta e coincide com o gabarito do Poliedro: são 10 PAs, com r $\in$ \{$\pm$1, $\pm$3, $\pm$5, $\pm$7, $\pm$9\} e a\_1 = (87-9r)/2, cujos termos vão de 3 a 84. Confirmei o resultado com o enunciado original e por força bruta. A resolução é até mais completa que o gabarito, porque justifica o limite |r| $\leq$ 9 e verifica os casos de fronteira. Não há divergência de mérito; no gabarito só aparecem pequenas falhas de digitação. \\
2 & 10,0 & A resolução está correta e é equivalente ao gabarito: S = (log$_{2}$(($\surd$5-1)/2), 1) = (-log$_{2}$((1+$\surd$5)/2), 1) $\approx$ (-0,694; 1). Primeiro ela mostra que log$_{2}$(2\^{}x+2) > 1 > 0, o que tira o módulo da esquerda. Depois resolve |2x| < log$_{2}$(2\^{}x+2) como duas desigualdades: a primeira dá x<1 e a segunda, via (t+1)(t$^2$+t-1)>0, dá x>log$_{2}$(($\surd$5-1)/2). Por fim confere que nos extremos vale a igualdade. Refiz as contas e fiz uma varredura numérica, e tudo confirma o intervalo. O gabarito só tem um erro de digitação sem consequência (< no lugar de > no Caso 2). O enunciado transcrito confere com o original. \\
3 & 10,0 & A resolução está correta e coincide com o gabarito do Poliedro nos dois itens: r = 1/16 e lado $\surd$(4 - 2$\surd$2)/2 $\approx$ 0,541. O método é o mesmo nos dois: z\^{}8 = ri, octógono regular de raio r\^{}\{1/8\}, área 2$\surd$2$\cdot$R$^2$ = $\surd$2, o que dá R = $\surd$2/2, e lei dos cossenos para o lado. Conferi numericamente as raízes, a área e o lado. A transcrição do enunciado confere com a prova original (página impressa 4, que é a página 3 do PDF). Não há divergência nem erro a apontar. \\
4 & 10,0 & A resolução e o gabarito concordam: a = -2 e p(x) = (x+1)$^2$(x-2)(x$^2$+1)(x$^2$-6x+1). As raízes são -1 (dupla), 2, $\pm$i e 3 $\pm$ 2$\surd$2. Refiz em Python, com aritmética exata, a divisão para a geral (o resto é (a+2)(5x$^2$+26x+13)), a divisão para a = -2, a fatoração e as multiplicidades, e está tudo certo. A resolução está correta e completa: mostra que as condições são necessárias e suficientes, verifica as multiplicidades e confere por Vieta. Não há erro a apontar em nenhum dos dois. \\
5 & 10,0 & A resolução está certa e confere com o gabarito do Poliedro: x\_C = 5$\surd$3 - 9 $\approx$ -0,340. As duas usam o mesmo método. Pela condição do perímetro, C fica numa segunda elipse com focos em A=(1,2) e B=(1,-2), com 2a=6 e b$^2$=5. A interseção com 2x$^2$+y$^2$=6 leva a x$^2$+18x+6=0, e a raiz -9-5$\surd$3 é descartada porque dá y$^2$ < 0. Refiz as contas em Python e o perímetro dá 10. O gabarito só tem erros de digitação, sem efeito no resultado. \\
6 & 10,0 & A resolução está correta e coincide com o gabarito do Poliedro: são 235 sêxtuplas. O escalonamento, o intervalo x6 $\in$ \{1,2,3\} e a contagem por estrelas e barras (91 + 78 + 66) estão certos. Confirmei o resultado por força bruta e pela contagem complementar 455 - 220. A resolução também justifica a equivalência dos sistemas e verifica a resposta com mais cuidado que o gabarito. Não há erro em nenhum dos dois. \\
7 & 10,0 & A resolução está correta e a resposta coincide com a do gabarito: n $\in$ \{-6, -4, 6, 8\}. Uma busca por força bruta em python e a divisão polinomial refeita confirmam esse conjunto e os quocientes -10, -30, -10 e 18. As únicas observações são de redação: uma provável perda de parênteses no discriminante do gabarito e uma frase pouco precisa no passo 5 da resolução. Nenhuma delas afeta a nota. \\
8 & 10,0 & A resolução e o gabarito do Poliedro estão certos e concordam nos dois itens. No item (a), x $\in$ (0,1{]} ? {[}9,10), porque o discriminante de t$^2$ - 6t + x(10-x) é proporcional a (x-1)(x-9). No item (b), a área 30 leva a (x-5)(y-3) = 0. Como x = 5 é impossível, y = 3 e x = 1 ou x = 9, e os lados medem $\surd$10 e 3$\surd$10. A resolução é mais completa que o gabarito: prova a recíproca e diz qual lado é qual em cada caso. Achei só um erro de digitação no gabarito ('x-5y-3 = 0' no lugar de (x-5)(y-3) = 0), sem efeito no resultado. O enunciado original é a página impressa 9, que corresponde à página 8 do PDF, e a transcrição confere com ele. \\
9 & 10,0 & A resolução está correta e chega ao mesmo resultado do gabarito: a área é 10+5$\surd$3 = 5(2+$\surd$3) $\approx$ 18,66. A seção é um trapézio de bases 8 e 2 e altura 2+$\surd$3, valor igual ao tg75$^\circ$ do gabarito. Conferi a altura da pirâmide, o plano da seção, os pontos médios de SE e SF e a área numericamente, e tudo bate. Explicar por que a seção tem de ser paralela a SBC ou SEF é um ponto a favor da resolução. O único problema está no texto do gabarito, em que a fórmula final aparece truncada, mas a resposta dele está certa. \\
10 & 10,0 & A resolução e o gabarito concordam: o único inteiro é k = 3. Refiz a conta e ela confirma isso. No 1º quadrante, f $\in$ ($\surd$2-1, 2-$\surd$2{]} ? (0,1), sem inteiros. No 3º quadrante, f $\in$ ($\surd$2+1, 2+$\surd$2{]} $\approx$ (2,41; 3,41{]}, e o 3 é atingido em x $\approx$ 3,429 e x $\approx$ 4,425. Uma amostragem numérica de f nos dois quadrantes confirmou esses intervalos. A resolução está correta, completa e mais rigorosa que o gabarito, que tem só pequenas imprecisões de notação sem efeito na resposta. Não consegui salvar uma nota de memória sobre a numeração das páginas do PDF porque a escrita foi bloqueada pela restrição de acesso. \\
\bottomrule\end{longtable}%
\section*{Questão 1}%
\href{https://poliedroresolve.sistemapoliedro.com.br/vestibulares/ita/2027/ita-2-fase-dia-1-20-10-2026-2-fase-dia-1/dissertativa-1-matematica-2-fase-dia-1-ita-2027}{Resolução do gabarito}\par\smallskip%
\begin{longtable}{@{}p{0.06\textwidth} p{0.33\textwidth} p{0.33\textwidth} p{0.1\textwidth} p{0.08\textwidth}@{}}\toprule Item & Resolução & Gabarito & Veredito & Certo \\ \midrule
Q1 & 10 PAs ordenadas, a\_k = (87 + (2k-11)r)/2, k = 1..10, com r $\in$ \{$\pm$1, $\pm$3, $\pm$5, $\pm$7, $\pm$9\} e a\_1 = (87 - 9r)/2. Tabela explícita: r=9: 3..84; r=7: 12..75; r=5: 21..66; r=3: 30..57; r=1: 39..48; e as 5 decrescentes correspondentes. Como conjuntos de números, são 5. & 10 PAs: crescentes (39,...,48) r=1; (30,...,57) r=3; (21,...,66) r=5; (12,...,75) r=7; (3,...,84) r=9; e as 5 decrescentes correspondentes (r = -1, -3, -5, -7, -9). & igual & ambos \\
\bottomrule\end{longtable}%
\begin{itemize}%
\item%
\textbf{Q1)} Conferi com o enunciado original (pág. 2 da prova, onde N = \{1,2,3,...\}). Da soma vem 2a\_1 + 9r = 87, então r é ímpar. Para todos os termos serem $\geq$ 1 é preciso 1 $\leq$ a\_1 $\leq$ 86, o que dá -9 $\leq$ r $\leq$ 9. Testei por força bruta em python (r de -20 a 20, a\_1 de 0 a 199, exigindo todos os termos $\geq$ 1): saem exatamente os 10 pares (r, a\_1) = ($\pm$1; 39/48), ($\pm$3; 30/57), ($\pm$5; 21/66), ($\pm$7; 12/75), ($\pm$9; 3/84). Os termos da tabela da resolução batem com os do gabarito, e a fórmula fechada a\_k = (87+(2k-11)r)/2 vale em k=1 e em k=10. Os dois chegam às mesmas 10 PAs: 5 crescentes e as 5 decrescentes correspondentes. A observação extra da resolução (são 5 conjuntos) não contradiz o gabarito.%
\end{itemize}%
\paragraph{Pontos de dúvida}%
\begin{itemize}%
\item%
\textit{(key)} A resolução do gabarito diz só que 'a razão deve ser ímpar' e lista r = 1, 3, 5, 7, 9. Não justifica por que para em r = 9 (r = 11 daria a\_1 = -6) nem trata as razões negativas pela inequação. A resolução avaliada faz isso direito (-9 $\leq$ r $\leq$ 9 a partir de 1 $\leq$ a\_1 $\leq$ 86).%
\item%
\textit{(key)} Erros de digitação na transcrição do gabarito: faltam parênteses em 'S\_10 = (a\_1 + a\_10 $\cdot$ 10)/2' e '(2a\_1 + 9r $\cdot$ 10)/2' (deveria ser (a\_1 + a\_10)$\cdot$10/2). As 5 PAs decrescentes aparecem grudadas numa sequência só ('...40,3957,54,...'). Nenhum dos dois muda a resposta.%
\item%
\textit{(statement)} 'Descreva todas as PAs' pode ser lido como contar sequências ordenadas (10) ou conjuntos de números (5). O gabarito conta 10, com as decrescentes. A resolução dá 10 como resposta principal e menciona os 5 conjuntos só como observação, então fica coberta nas duas leituras. Ela também verifica que contar o 0 como natural não muda o resultado.%
\end{itemize}%
\section*{Questão 2}%
\href{https://poliedroresolve.sistemapoliedro.com.br/vestibulares/ita/2027/ita-2-fase-dia-1-20-10-2026-2-fase-dia-1/dissertativa-2-matematica-2-fase-dia-1-ita-2027}{Resolução do gabarito}\par\smallskip%
\begin{longtable}{@{}p{0.06\textwidth} p{0.33\textwidth} p{0.33\textwidth} p{0.1\textwidth} p{0.08\textwidth}@{}}\toprule Item & Resolução & Gabarito & Veredito & Certo \\ \midrule
Questão 2 (conjunto solução) & S = (-log$_{2}$((1+$\surd$5)/2), 1) = (log$_{2}$(($\surd$5-1)/2), 1) $\approx$ (-0,694; 1) & S = \{ x $\in$ $\mathbb{R}$ | log$_{2}$(($\surd$5 - 1)/2) < x < 1 \} & equivalente & ambos \\
\bottomrule\end{longtable}%
\begin{itemize}%
\item%
\textbf{Questão 2 (conjunto solução))} É o mesmo intervalo aberto. A resolução também escreve o extremo inferior como -log$_{2}$$\varphi$, com $\varphi$=(1+$\surd$5)/2; isso é igual a log$_{2}$(($\surd$5-1)/2), porque ($\surd$5-1)/2 = 1/$\varphi$. Conferi numericamente: os dois valores dão -0,694242. Refiz a conta: log$_{2}$(2\^{}x+2) > 1 para todo x real, então o módulo da esquerda some. A condição log$_{2}$(t+2) > 2x, com t=2\^{}x, vira t$^2$-t-2<0 e dá x<1. A condição log$_{2}$(t+2) > -2x vira t$^3$+2t$^2$-1 = (t+1)(t$^2$+t-1) > 0 e dá t > ($\surd$5-1)/2. Nos dois extremos vale a igualdade, então ambos ficam excluídos. Uma varredura com 1 200 001 pontos em {[}-6, 6{]} encontrou a desigualdade valendo só em (-0,69424; 0,99999), o que confirma o resultado. O gabarito separa os casos x$\geq$0 e x<0. A resolução usa |a|<L $\Leftrightarrow$ -L<a<L e faz a interseção das duas condições. Os dois caminhos estão corretos.%
\end{itemize}%
\paragraph{Pontos de dúvida}%
\begin{itemize}%
\item%
\textit{(key)} Erro de digitação no Caso 2 do gabarito: aparece "log$_{2}$(2\^{}x + 2) < -2x", mas o certo é ">". A linha seguinte (2\^{}x + 2 > 2\^{}(-2x)) já usa o sentido correto, então o resultado final não muda.%
\item%
\textit{(solution)} No passo 3 da resolução aparece "para todo k real inteiro", uma redação imprecisa (basta k inteiro, e a identidade 2\^{}\{kx\}=t\^{}k vale até para k real). Isso não afeta a solução.%
\item%
\textit{(solution)} A "conferência numérica adicional" com 200 000 pontos é citada, mas o resultado não aparece. Ela não faz parte da demonstração, e a refiz aqui: o resultado se confirma.%
\end{itemize}%
\section*{Questão 3}%
\href{https://poliedroresolve.sistemapoliedro.com.br/vestibulares/ita/2027/ita-2-fase-dia-1-20-10-2026-2-fase-dia-1/dissertativa-3-matematica-2-fase-dia-1-ita-2027}{Resolução do gabarito}\par\smallskip%
\begin{longtable}{@{}p{0.06\textwidth} p{0.33\textwidth} p{0.33\textwidth} p{0.1\textwidth} p{0.08\textwidth}@{}}\toprule Item & Resolução & Gabarito & Veredito & Certo \\ \midrule
r & r = 1/16 = 0,0625 & r = 1/16 & igual & ambos \\
comprimento dos lados & s = $\surd$(1 - $\surd$2/2) = $\surd$(4 - 2$\surd$2)/2 $\approx$ 0,541 (os oito lados são iguais) & l = $\surd$(4 - 2$\surd$2)/2 = $\surd$((2 - $\surd$2)/2) & igual & ambos \\
\bottomrule\end{longtable}%
\begin{itemize}%
\item%
\textbf{r)} De z\^{}8 = ri = r$\cdot$cis($\pi$/2), as raízes são vértices de um octógono regular de raio R = r\^{}\{1/8\}. A área é 8$\cdot$(1/2)R$^2$$\cdot$sen45$^\circ$ = 2$\surd$2$\cdot$R$^2$. Igualando a $\surd$2, R$^2$ = 1/2, logo r = R\^{}8 = (1/2)\^{}4 = 1/16. Conferi numericamente: com r = 1/16, |r + i$\cdot$z\_k\^{}8| $\approx$ 1e{-}16 para os oito z\_k, e a área pela fórmula do laço (shoelace) dá 1,41421\ldots{} = $\surd$2.%
\item%
\textbf{comprimento dos lados)} Pela lei dos cossenos com R$^2$ = 1/2: l$^2$ = 1/2 + 1/2 - 2$\cdot$(1/2)$\cdot$cos45$^\circ$ = 1 - $\surd$2/2, então l = $\surd$(4 - 2$\surd$2)/2. A resolução escreve essa mesma forma, e $\surd$(1 - $\surd$2/2) e $\surd$((2 - $\surd$2)/2) são só outras grafias do mesmo número. A conferência numérica bate: |z\_1 - z\_0| = 0,541196\ldots{}, igual às três expressões.%
\end{itemize}%
\paragraph{Pontos de dúvida}%
\begin{itemize}%
\item%
\textit{(statement)} O enunciado não diz em que ordem os vértices são ligados. A resolução e o gabarito tomam o octógono convexo, com os vértices ligados na ordem dos argumentos, que é a leitura natural. A resolução justifica isso com o ângulo central de $\pi$/4 entre raízes consecutivas.%
\end{itemize}%
\section*{Questão 4}%
\href{https://poliedroresolve.sistemapoliedro.com.br/vestibulares/ita/2027/ita-2-fase-dia-1-20-10-2026-2-fase-dia-1/dissertativa-4-matematica-2-fase-dia-1-ita-2027}{Resolução do gabarito}\par\smallskip%
\begin{longtable}{@{}p{0.06\textwidth} p{0.33\textwidth} p{0.33\textwidth} p{0.1\textwidth} p{0.08\textwidth}@{}}\toprule Item & Resolução & Gabarito & Veredito & Certo \\ \midrule
Questão 4 (valor de a e raízes de p) & a = -2; p(x) = (x+1)\^{}2(x-2)(x\^{}2+1)(x\^{}2-6x+1); raízes: -1 (dupla), 2, i, -i, 3+2$\surd$2, 3-2$\surd$2 & a = -2; raízes: -1 (dupla), 2, i, -i, 3+2$\surd$2, 3-2$\surd$2 & igual & ambos \\
\bottomrule\end{longtable}%
\begin{itemize}%
\item%
\textbf{Questão 4 (valor de a e raízes de p))} Conferi tudo de novo em Python, com aritmética exata. Para a geral, o resto da divisão de p por q é (a+2)(5x$^2$+26x+13). Ele só se anula para a = -2, que é o único valor possível. Também conferi p(2) = 85(a+2), p(-1) = -8(a+2) e p'(-1) = 16(a+2), que batem com o passo 3 da resolução. Para a = -2, o quociente dá x$^{4}$-6x$^3$+2x$^2$-6x+1 com resto 0, e o produto (x+1)$^2$(x-2)(x$^2$+1)(x$^2$-6x+1) reproduz exatamente os coeficientes {[}1,-6,-1,10,7,14,9,-2{]}. Como s(-1) = 16 $\neq$ 0 e s(2) = -35 $\neq$ 0, as multiplicidades são exatamente 2 para -1 e 1 para 2. As contas de Vieta da resolução estão certas: e1 = 6, e2 = -1 e e7 = 2. A resolução e o gabarito dão a mesma resposta.%
\end{itemize}%
\paragraph{Pontos de dúvida}%
\begin{itemize}%
\item%
\textit{(statement)} A transcrição indica 'página 5', mas a Questão 04 está na página 4 do PDF, que tem o número 5 impresso. A página 5 do PDF traz a Questão 05. A transcrição confere com o original.%
\item%
\textit{(key)} O gabarito acha a usando só p(-1) = 0 (-8a-16 = 0, que ele escreve como 8a+16 = 0). Não usa p(2) = 0 nem p'(-1) = 0, porque a raiz -1 é dupla. A conclusão é válida, porque a divisão explícita feita depois tem resto nulo. A resolução é mais completa: impõe as três condições e confirma que bastam.%
\end{itemize}%
\section*{Questão 5}%
\href{https://poliedroresolve.sistemapoliedro.com.br/vestibulares/ita/2027/ita-2-fase-dia-1-20-10-2026-2-fase-dia-1/dissertativa-5-matematica-2-fase-dia-1-ita-2027}{Resolução do gabarito}\par\smallskip%
\begin{longtable}{@{}p{0.06\textwidth} p{0.33\textwidth} p{0.33\textwidth} p{0.1\textwidth} p{0.08\textwidth}@{}}\toprule Item & Resolução & Gabarito & Veredito & Certo \\ \midrule
abscissa de C & x\_C = {-}9 + 5$\surd$3 $\approx$ {-}0,340 (dois vértices possíveis, (-9+5$\surd$3, $\pm$3$\surd$(20$\surd$3-34)), com a mesma abscissa) & x = 5$\surd$3 - 9 & igual & ambos \\
\bottomrule\end{longtable}%
\begin{itemize}%
\item%
\textbf{abscissa de C)} O valor é o mesmo, só escrito em outra ordem. Refiz a conta: com A=(1,2) e B=(1,-2), AB=4 e CA+CB=6. Então C está na elipse de focos A e B, (x-1)$^2$/5 + y$^2$/9 = 1, ou seja, 9x$^2$ - 18x + 5y$^2$ = 36. Fazendo 5$\cdot$(2x$^2$+y$^2$=6) menos essa equação, sai x$^2$ + 18x + 6 = 0, que dá x = -9 $\pm$ 5$\surd$3. A raiz -9-5$\surd$3 $\approx$ -17,66 está fora de |x| $\leq$ $\surd$3, pois dá y$^2$ < 0, e por isso é descartada. Conferi em Python: x $\approx$ -0,33975, y $\approx$ $\pm$2,40190, 2x$^2$+y$^2$ = 6, CA $\approx$ 1,39873, CB $\approx$ 4,60127 e perímetro = 10,0000. As duas resoluções descartam a raiz espúria corretamente. A resolução avaliada mostra ainda que o triângulo não é degenerado (x\_C $\neq$ 1).%
\end{itemize}%
\paragraph{Pontos de dúvida}%
\begin{itemize}%
\item%
\textit{(key)} O gabarito tem erros de digitação: escreve "(x{-}1\^{}2)/5" no lugar de "(x{-}1)\^{}2/5" e "A = 1,2" sem parênteses. A equação desenvolvida, 9x$^2$-18x+5y$^2$=36, está certa, então esses erros não afetam o resultado.%
\item%
\textit{(solution)} A resolução usa A$\neq$B, o que é necessário para haver triângulo, e escolhe A=(1,2) e B=(1,-2) sem perda de generalidade. As duas hipóteses são razoáveis e o resultado não depende da escolha.%
\end{itemize}%
\section*{Questão 6}%
\href{https://poliedroresolve.sistemapoliedro.com.br/vestibulares/ita/2027/ita-2-fase-dia-1-20-10-2026-2-fase-dia-1/dissertativa-6-matematica-2-fase-dia-1-ita-2027}{Resolução do gabarito}\par\smallskip%
\begin{longtable}{@{}p{0.06\textwidth} p{0.33\textwidth} p{0.33\textwidth} p{0.1\textwidth} p{0.08\textwidth}@{}}\toprule Item & Resolução & Gabarito & Veredito & Certo \\ \midrule
Questão 6 (número de sêxtuplas) & 235 sêxtuplas & 235 sêxtuplas & igual & ambos \\
\bottomrule\end{longtable}%
\begin{itemize}%
\item%
\textbf{Questão 6 (número de sêxtuplas))} Conferi o enunciado no original (página impressa 7, que é a página 6 do PDF), e ele bate com a transcrição. O escalonamento dá x2 = 6, x1 = 8 - 2x6 e x3 + x4 + x5 = 16 - x6. Como $\mathbb{N}$ = \{1,2,3,...\}, x6 só pode ser 1, 2 ou 3. As contagens são C(14,2) + C(13,2) + C(12,2) = 91 + 78 + 66 = 235. Uma busca exaustiva em Python, com todas as variáveis $\geq$ 1, também encontrou 235. A contagem complementar da resolução, C(15,3) - C(12,3) = 455 - 220 = 235, está certa. A observação de que daria 605 se o 0 fosse natural também foi confirmada por força bruta. A resolução e o gabarito seguem o mesmo caminho. A notação P\_14\^{}\{12,2\} = 14!/(12!$\cdot$2!) = 91 do gabarito equivale a C(14,2).%
\end{itemize}%
\paragraph{Pontos de dúvida}%
\begin{itemize}%
\item%
\textit{(key)} O gabarito afirma 1 $\leq$ x6 $\leq$ 3 sem justificar. Isso vem de x1 = 8 - 2x6 $\geq$ 1 e de x6 $\geq$ 1. Falta só redação: o resultado está certo.%
\item%
\textit{(statement)} A resposta depende da convenção $\mathbb{N}$ = \{1,2,3,...\} da página 2 da prova. Se o 0 fosse natural, seriam 605 sêxtuplas (confirmado por força bruta). Como a prova fixa a convenção, não há ambiguidade de fato.%
\end{itemize}%
\section*{Questão 7}%
\href{https://poliedroresolve.sistemapoliedro.com.br/vestibulares/ita/2027/ita-2-fase-dia-1-20-10-2026-2-fase-dia-1/dissertativa-7-matematica-2-fase-dia-1-ita-2027}{Resolução do gabarito}\par\smallskip%
\begin{longtable}{@{}p{0.06\textwidth} p{0.33\textwidth} p{0.33\textwidth} p{0.1\textwidth} p{0.08\textwidth}@{}}\toprule Item & Resolução & Gabarito & Veredito & Certo \\ \midrule
Questão 7 & n $\in$ \{-6, -4, 6, 8\} (quocientes -10, -30, -10 e 18) & n $\in$ \{-6, -4, 6, 8\} & igual & ambos \\
\bottomrule\end{longtable}%
\begin{itemize}%
\item%
\textbf{Questão 7)} Conferi o enunciado no original (página impressa 8, que é a página 7 do PDF) e ele bate com a transcrição. Refiz as contas: 2n$^3$-5n$^2$-52n+90 = (2n-1)(n$^2$-2n-27) + 63, então o quociente é inteiro se e somente se n$^2$-2n-27 divide 63. Como n$^2$-2n-27 = (n-1)$^2$ - 28, é preciso que d+28 seja quadrado perfeito, sendo d um divisor de 63. Isso só ocorre para d = -3 (25), que dá n = 6 ou n = -4, e para d = 21 (49), que dá n = 8 ou n = -6. Uma busca por força bruta em python, com n de -100000 a 100000, encontrou exatamente \{-6, -4, 6, 8\}, com quocientes -10, -30, -10 e 18. Esses valores coincidem com os da resolução. Os dois métodos são equivalentes: a resolução completa o quadrado e o gabarito usa o discriminante.%
\end{itemize}%
\paragraph{Pontos de dúvida}%
\begin{itemize}%
\item%
\textit{(key)} Na resolução do gabarito, como transcrita, o discriminante aparece como '$\Delta$ = 4 + 4k + 27 = m$^2$, m$^2$ = 4k + 28'. O correto é $\Delta$ = 4 + 4(k+27) = 4(k+28). Parece perda de parênteses na transcrição. A conclusão de que k+28 precisa ser quadrado perfeito está certa e o resultado final não muda.%
\item%
\textit{(solution)} No passo 5, a resolução diz que 'em cada caso o número está estritamente entre dois quadrados consecutivos', mas a lista inclui -35, que é negativo. Esse caso já tinha sido descartado à parte por ser negativo, então é só uma imprecisão de redação, sem efeito no resultado.%
\end{itemize}%
\section*{Questão 8}%
\href{https://poliedroresolve.sistemapoliedro.com.br/vestibulares/ita/2027/ita-2-fase-dia-1-20-10-2026-2-fase-dia-1/dissertativa-8-matematica-2-fase-dia-1-ita-2027}{Resolução do gabarito}\par\smallskip%
\begin{longtable}{@{}p{0.06\textwidth} p{0.33\textwidth} p{0.33\textwidth} p{0.1\textwidth} p{0.08\textwidth}@{}}\toprule Item & Resolução & Gabarito & Veredito & Certo \\ \midrule
a & 0 < x $\leq$ 1 ou 9 $\leq$ x < 10, isto é, x $\in$ (0,1{]} ? {[}9,10) & 0 < x $\leq$ 1 ou 9 $\leq$ x < 10 & igual & ambos \\
b & Lados $\surd$10 e 3$\surd$10 ($\approx$ 3,16 e 9,49). Ocorre só para x = 1 ou x = 9, com BF = AH = 3. & Lados 3$\surd$10 e $\surd$10 & igual & ambos \\
\bottomrule\end{longtable}%
\begin{itemize}%
\item%
\textbf{a)} As duas respostas são iguais. A condição y(6-y) = x(10-x), com BF = y e AH = 6-y, é a mesma nas duas resoluções: a resolução a obtém pelo produto escalar EF$\cdot$EH = 0 e o gabarito pela semelhança entre os triângulos EBF e HAE. O discriminante 36 - 4x(10-x) $\geq$ 0 equivale a (x-1)(x-9) $\geq$ 0. Conferi numericamente: $\Delta$ = 0 em x = 1 e x = 9, $\Delta$ > 0 em x = 0,5 e x = 9,5, e $\Delta$ < 0 em x = 3 e x = 5. A resolução ainda prova a recíproca (toda raiz real gera um retângulo com G e H no interior dos lados) e mostra que 0 < t < 6, o que deixa o intervalo bem justificado.%
\item%
\textbf{b)} As duas respostas são iguais. Refiz a conta: S = 60 - (10-x)y - x(6-y) = 30 leva a 2(x-5)(y-3) = 0. O caso x = 5 é impossível, porque daria yh = 25 > 9. Sobra y = 3, e então x(10-x) = 9 dá x = 1 ou x = 9. Para (x, y) = (1, 3), EF = $\surd$(81+9) = 3$\surd$10 $\approx$ 9,487 e EH = $\surd$(1+9) = $\surd$10 $\approx$ 3,162. Para x = 9 os valores se trocam. A área fica 3$\surd$10$\cdot$$\surd$10 = 30. A resolução chega à área pela fórmula |det(EF, EH)|, que confere com 60 menos os quatro triângulos dos cantos.%
\end{itemize}%
\paragraph{Pontos de dúvida}%
\begin{itemize}%
\item%
\textit{(key)} O gabarito tem um erro de digitação no item (b). Escreve 'xy-3x-5y+15 = 0 $\Rightarrow$ x-5y-3 = 0', mas o certo é (x-5)(y-3) = 0. O argumento seguinte ('x = 5 é impossível, logo y = 3') mostra que a fatoração pretendida era essa, então o erro não altera a resposta.%
\item%
\textit{(statement)} As duas resoluções supõem que 'inscrito' quer dizer um vértice de EFGH em cada lado de ABCD, na ordem cíclica (G em CD e H em DA). É a leitura usual e o enunciado não a declara. A resolução a conferir deixa a hipótese explícita.%
\end{itemize}%
\section*{Questão 9}%
\href{https://poliedroresolve.sistemapoliedro.com.br/vestibulares/ita/2027/ita-2-fase-dia-1-20-10-2026-2-fase-dia-1/dissertativa-9-matematica-2-fase-dia-1-ita-2027}{Resolução do gabarito}\par\smallskip%
\begin{longtable}{@{}p{0.06\textwidth} p{0.33\textwidth} p{0.33\textwidth} p{0.1\textwidth} p{0.08\textwidth}@{}}\toprule Item & Resolução & Gabarito & Veredito & Certo \\ \midrule
área da seção & 10 + 5$\surd$3 = 5(2+$\surd$3) $\approx$ 18,66 u.a. (trapézio ADXY de bases 8 e 2 e altura 2+$\surd$3) & 5(2 + $\surd$3) & equivalente & ambos \\
\bottomrule\end{longtable}%
\begin{itemize}%
\item%
\textbf{área da seção)} 10+5$\surd$3 é o mesmo número que 5(2+$\surd$3), escrito de outro jeito. Refiz a conta: com AB=4 e ângulo do ápice $\pi$/6, a lei dos cossenos dá L$^2$=16(2+$\surd$3), logo H=4$\surd$(1+$\surd$3)$\approx$6,6116. Para a seção passar por AD, ela tem de ser paralela a SBC ou a SEF; as duas escolhas dão a mesma área, por simetria. O plano corta SE e SF nos pontos médios, então a seção é um trapézio de bases 8 e 2. A distância entre as bases é $\surd$(3+H$^2$/4)=$\surd$(7+4$\surd$3)=2+$\surd$3, o mesmo valor que o gabarito chama de OT=tg75$^\circ$. Também conferi o apótema lateral: $\surd$(H$^2$+12)=2(2+$\surd$3)=2$\cdot$tg75$^\circ$$\approx$7,4641, igual ao VM do gabarito. Calculei a área do quadrilátero A,Y,X,D por produto vetorial e obtive 18,6603, que bate com 10+5$\surd$3. Também verifiquei que a normal desse plano é paralela à normal da face SBC.%
\end{itemize}%
\paragraph{Pontos de dúvida}%
\begin{itemize}%
\item%
\textit{(key)} O texto da resolução do gabarito parece ter perdido símbolos na transcrição. A linha final aparece como "(AD + PQOT)/2 = (8 + 2 $\cdot$ tg75$^\circ$)/2 = 52 + $\surd$(3)". O certo seria (AD+PQ)$\cdot$OT/2 = (8+2)$\cdot$tg75$^\circ$/2 = 5(2+$\surd$3). A resposta por item (5(2+$\surd$3)) está correta e coincide com a da resolução, então isso não altera a nota.%
\item%
\textit{(statement)} O enunciado diz "paralela a uma das faces laterais" sem dizer qual. Só as faces opostas aos lados paralelos à diagonal (SBC e SEF) permitem um plano paralelo que contenha a diagonal AD, e as duas dão a mesma área por simetria. A resolução justifica isso no passo 3. O gabarito escolhe a face sem justificar.%
\item%
\textit{(statement)} A indicação "páginas 10" corresponde à página impressa 10 da prova, que é a página 9 do PDF. A página 10 do PDF traz a Questão 10. Comparei o enunciado da Questão 09 com o original e ele está igual à transcrição.%
\end{itemize}%
\section*{Questão 10}%
\href{https://poliedroresolve.sistemapoliedro.com.br/vestibulares/ita/2027/ita-2-fase-dia-1-20-10-2026-2-fase-dia-1/dissertativa-10-matematica-2-fase-dia-1-ita-2027}{Resolução do gabarito}\par\smallskip%
\begin{longtable}{@{}p{0.06\textwidth} p{0.33\textwidth} p{0.33\textwidth} p{0.1\textwidth} p{0.08\textwidth}@{}}\toprule Item & Resolução & Gabarito & Veredito & Certo \\ \midrule
Q10 --- inteiros k tais que f(x)=k tem solução no 1º ou 3º quadrante & k = 3 (único). No 1º quadrante, f $\in$ ($\surd$2-1, 2-$\surd$2{]} ? (0,1). No 3º quadrante, f $\in$ ($\surd$2+1, 2+$\surd$2{]} ? (2,4), e 3 é atingido em x = 3$\pi$/4 + arcsen(3 - 3$\surd$2/2) $\approx$ 3,429 e em x = 7$\pi$/4 - arcsen(3 - 3$\surd$2/2) $\approx$ 4,425. & k = 3 & igual & ambos \\
\bottomrule\end{longtable}%
\begin{itemize}%
\item%
\textbf{Q10 --- inteiros k tais que f(x)=k tem solução no 1º ou 3º quadrante)} O enunciado confere com o original (página 10 do PDF, que traz o número impresso 11). Refiz as contas: f(x) = $\surd$2(|cos x|+|sen x|) - (sen x + cos x). No 1º quadrante, f = ($\surd$2-1)(sen x + cos x) = (2-$\surd$2)$\cdot$sen(x+$\pi$/4), com valores em ($\surd$2-1, 2-$\surd$2{]} $\approx$ (0,414; 0,586{]}. No 3º quadrante, f = -($\surd$2+1)(sen x + cos x) = -(2+$\surd$2)$\cdot$sen(x+$\pi$/4), com valores em ($\surd$2+1, 2+$\surd$2{]} $\approx$ (2,414; 3,414{]}. Conferi numericamente com python, amostrando f em 400 mil pontos de cada quadrante: os extremos batem. Também avaliei f nas duas soluções explícitas e obtive f = 3,000. A resposta k = 3 está correta e coincide com o gabarito.%
\end{itemize}%
\paragraph{Pontos de dúvida}%
\begin{itemize}%
\item%
\textit{(key)} O gabarito escreve f(x) = $\surd$2 cos x + $\surd$2 sen x - sen x - cos x de forma geral, sem módulos, e só depois corrige o sinal caso a caso. Além disso, toma intervalos fechados (0 $\leq$ x $\leq$ $\pi$/2 e $\pi$ $\leq$ x $\leq$ 3$\pi$/2), o que inclui os eixos, e mistura desigualdades estritas e não estritas ($\surd$2/2 < \ldots{} depois $\surd$2-1 $\leq$ f). Isso não muda a resposta, porque os valores de f nos eixos ($\surd$2-1 e $\surd$2+1) não são inteiros.%
\item%
\textit{(key)} O gabarito não justifica explicitamente que o valor 3 é atingido. Isso fica implícito no fato de sen(x+$\pi$/4) percorrer todo o intervalo. A resolução conferida trata esse ponto com mais rigor: usa o teorema do valor intermediário e exibe as soluções explícitas.%
\item%
\textit{(solution)} Imprecisões de redação na resolução, sem efeito no resultado: (i) no Passo 5, o cálculo de f(5$\pi$/4) usa cos(5$\pi$/2) = 0, o que está certo porque 2x = 5$\pi$/2, mas a frase pode confundir quem lê; (ii) a desigualdade final do Passo 3 encadeia 0 < $\surd$2-1 < f, o que está correto.%
\end{itemize}%
\end{document}